\documentclass[11pt,a4paper]{article}
\usepackage[utf8]{inputenc}
\usepackage[T1]{fontenc}
\usepackage[french]{babel}
\usepackage[top=2cm,bottom=2cm,left=2cm,right=2cm]{geometry}
\usepackage{xcolor}
\usepackage{tcolorbox}
\tcbuselibrary{skins,breakable}
\usepackage{fancyhdr}
\usepackage{enumitem}
\usepackage{lmodern}

% --- Couleur discipline : DROIT ---
\definecolor{mainColor}{RGB}{160,60,0}
\definecolor{darkGray}{RGB}{50,50,50}

\tcbset{boxrule=0.5pt, arc=3pt, left=5pt, right=5pt, top=3pt, bottom=3pt, breakable}

\newtcolorbox{defbox}[1]{
  colback=mainColor!8, colframe=mainColor, coltitle=white,
  fonttitle=\bfseries\footnotesize, title=DÉFINITION \textemdash\ #1}

\newtcolorbox{keybox}{
  colback=darkGray!7, colframe=darkGray!50, coltitle=white,
  fonttitle=\bfseries\footnotesize, title=POINT CLÉ}

\newtcolorbox{exbox}{
  colback=mainColor!5, colframe=mainColor!40,
  fonttitle=\bfseries\footnotesize, title=EXEMPLE}

\newtcolorbox{alertbox}{
  colback=red!6, colframe=red!70, coltitle=white,
  fonttitle=\bfseries\footnotesize, title=ATTENTION}

\pagestyle{fancy}\fancyhf{}
\fancyhead[L]{\small\textcolor{mainColor}{\textbf{CEJM — BTS SIO}}}
\fancyhead[C]{\small\textbf{Chapitre 3 — La formation du contrat}}
\fancyhead[R]{\small\textcolor{mainColor}{\textbf{DROIT}}}
\fancyfoot[C]{\small\thepage}
\renewcommand{\headrulewidth}{0.5pt}
\setlength{\headheight}{15pt}
\setlist{itemsep=2pt, topsep=3pt, parsep=0pt}

\begin{document}

\begin{center}
  {\Large\bfseries\textcolor{mainColor}{Chapitre 3}}\\[2pt]
  {\large La formation du contrat}\\[4pt]
  \textit{\textcolor{darkGray}{Quelles sont les conditions de validité d'un contrat et comment se forme-t-il ?}}
\end{center}
\vspace{4pt}\hrule\vspace{8pt}

\section*{Définitions essentielles}

\begin{defbox}{Contrat}
Accord de volontés entre deux ou plusieurs personnes destiné à créer, modifier ou éteindre des \textbf{obligations juridiques}.
\end{defbox}

\vspace{4pt}
\begin{defbox}{Offre (pollicitation)}
Proposition \textbf{ferme, précise et complète} de contracter, faite à une personne déterminée ou au public.
\end{defbox}

\vspace{4pt}
\begin{defbox}{Consentement}
Accord de volonté \textbf{libre et éclairé} de chaque partie au contrat.
\end{defbox}

\section*{Les 4 conditions de validité du contrat}

\begin{keybox}
\textbf{Consentement} + \textbf{Capacité} + \textbf{Contenu licite et certain} (objet + cause) → contrat valide.
\end{keybox}

\subsection*{1. Le consentement — libre et éclairé}
\begin{itemize}
  \item Doit être \textbf{libre} (sans contrainte) et \textbf{éclairé} (en connaissance de cause)
\end{itemize}

\textbf{Vices du consentement} (rendent le contrat annulable) :
\begin{itemize}
  \item \textbf{Erreur} : méconnaissance de la réalité (sur la substance de la chose)
  \item \textbf{Dol} : tromperie intentionnelle d'une partie pour obtenir l'accord
  \item \textbf{Violence} : contrainte physique ou morale sur le cocontractant
\end{itemize}

\subsection*{2. La capacité juridique}
\begin{itemize}
  \item \textbf{Capacité de jouissance} : aptitude à avoir des droits (toute personne)
  \item \textbf{Capacité d'exercice} : aptitude à exercer ses droits seul
  \item \textbf{Incapables} : mineurs non émancipés, majeurs sous tutelle/curatelle
\end{itemize}

\subsection*{3. Un contenu licite et certain}
\begin{itemize}
  \item L'objet doit être \textbf{déterminé ou déterminable}, \textbf{possible} et \textbf{licite}
  \item Conforme à la loi et aux bonnes mœurs (ex. : vente d'organes → nul)
\end{itemize}

\subsection*{4. La forme}
\begin{itemize}
  \item \textbf{Contrat consensuel} : accord de volonté suffit (la majorité des contrats)
  \item \textbf{Contrat solennel} : forme particulière exigée (ex. : acte notarié pour l'immobilier)
  \item \textbf{Contrat réel} : remise de la chose nécessaire (ex. : prêt)
\end{itemize}

\section*{Le processus de formation}

\begin{center}
\textbf{Pourparlers} $\rightarrow$ \textbf{Offre} $\rightarrow$ \textbf{Acceptation} $\rightarrow$ \textbf{Contrat formé}
\end{center}

\begin{itemize}
  \item L'\textbf{offre} : ferme + précise + complète (tous les éléments essentiels)
  \item L'\textbf{acceptation} : pure et simple (toute modification = contre-offre)
  \item Le contrat est formé dès la \textbf{rencontre de l'offre et de l'acceptation}
\end{itemize}

\begin{exbox}
\textbf{Vente en ligne} : offre = fiche produit sur le site, acceptation = validation de la commande, formation au moment du paiement.\\
\textbf{Achat en magasin} : offre = étiquetage, acceptation = passage en caisse.
\end{exbox}

\section*{Classification des contrats}

\begin{center}
\begin{tabular}{|l|l|}
\hline
\textbf{Type} & \textbf{Caractéristique} \\
\hline
Synallagmatique & Obligations \textbf{réciproques} (vente) \\
Unilatéral & Obligations d'\textbf{une seule partie} (donation) \\
À titre onéreux & \textbf{Avantages réciproques} \\
À titre gratuit & \textbf{Libéralité} d'une partie \\
De gré à gré & Négociation \textbf{libre} \\
D'adhésion & Conditions \textbf{imposées} (contrat type) \\
\hline
\end{tabular}
\end{center}

\section*{Les sanctions : la nullité}

\begin{alertbox}
La nullité \textbf{anéantit rétroactivement} le contrat : les parties retrouvent leur situation initiale (restitutions réciproques).
\end{alertbox}

\begin{itemize}
  \item \textbf{Nullité absolue} : violation de l'ordre public (ex. : objet illicite) — tout intéressé peut agir
  \item \textbf{Nullité relative} : protection d'une partie (ex. : vice du consentement) — délai : \textbf{5 ans}
\end{itemize}

\section*{Points clés à retenir}

\begin{keybox}
Un contrat valide nécessite : \textbf{consentement libre et éclairé}, \textbf{capacité}, \textbf{contenu licite et certain}.
\end{keybox}
\vspace{3pt}
\begin{keybox}
L'acceptation doit être \textbf{pure et simple} : toute modification constitue une contre-offre.
\end{keybox}
\vspace{3pt}
\begin{keybox}
En cas de vice du consentement (erreur, dol, violence), le contrat peut être \textbf{annulé} à la demande de la victime.
\end{keybox}

\section*{Mots-clés}
\textcolor{mainColor}{\textbf{Contrat}} $\cdot$
\textcolor{mainColor}{\textbf{Offre / Pollicitation}} $\cdot$
\textcolor{mainColor}{\textbf{Acceptation}} $\cdot$
\textcolor{mainColor}{\textbf{Consentement}} $\cdot$
\textcolor{mainColor}{\textbf{Vices du consentement}} $\cdot$
\textcolor{mainColor}{\textbf{Capacité juridique}} $\cdot$
\textcolor{mainColor}{\textbf{Objet licite}} $\cdot$
\textcolor{mainColor}{\textbf{Nullité absolue / relative}} $\cdot$
\textcolor{mainColor}{\textbf{Contrat synallagmatique}} $\cdot$
\textcolor{mainColor}{\textbf{Contrat d'adhésion}}

\end{document}
